\section{Economic Quantification}
\label{sec:economic}

\subsection{The Aggregate Efficiency Gain}

The US trucking industry generates approximately \$900 billion in annual revenue \citep{BLS_trucking2023}. At a 35\% empty-mile rate, carriers incur \$315 billion in zero-revenue operating costs annually---the fully-loaded cost of moving empty trucks. A 15-percentage-point reduction to 20\% empty reduces the zero-revenue cost base to \$180 billion, a savings of \$135 billion per year. The 13-percentage-point reduction to 22\% (the conservative bound) produces savings of \$117 billion per year.

These figures require careful interpretation. The \$135 billion is not a direct transfer to carrier profits; it is an efficiency gain that will be competed away through shipper rate negotiation, carrier cost reduction, driver compensation improvement, and market-rate adjustment. The economic analysis focuses on aggregate welfare gain---the reduction in total resources consumed to move a given quantity of freight---rather than the distributional question of who captures it.

\subsubsection{Operating Cost Basis}

Empty-mile costs are incurred at the truck operating cost rate, not the revenue rate. ATRI reports the following cost components per mile for a combination truck \citep{ATRI_costs2024}:

\begin{table}[h]
\centering
\caption{Truck Operating Costs per Mile (ATRI 2024)}
\label{tab:costs}
\begin{tabular}{lcc}
\toprule
Cost Component & \$/Mile & \% of Total \\
\midrule
Driver wages \& benefits & \$0.617 & 38.5\% \\
Fuel \& fuel taxes & \$0.457 & 28.5\% \\
Truck/trailer lease or purchase & \$0.196 & 12.2\% \\
Insurance & \$0.106 & 6.6\% \\
Maintenance \& repair & \$0.148 & 9.2\% \\
Permits \& licenses & \$0.022 & 1.4\% \\
Tolls & \$0.063 & 3.9\% \\
\midrule
Total & \$1.609 & 100\% \\
\bottomrule
\end{tabular}
\end{table}

At \$1.609 per mile (all-in empty-mile cost), the national empty-mile total is:

\begin{equation}
  C_{\text{empty}} = E_{\text{miles}} \times \$1.609
\end{equation}

Estimating total truck-miles: BTS data indicate approximately 300 billion truck-miles annually on US roads \citep{BLS_trucking2023}. At 35\% empty, the empty-mile total is approximately 105 billion miles per year. At \$1.609/mile:

\begin{equation}
  C_{\text{empty}} = 105 \times 10^9 \times \$1.609 = \$168.9 \text{ billion/year}
\end{equation}

A 15 percentage-point reduction (from 35\% to 20\%) reduces empty miles by 45 billion miles annually, saving:

\begin{equation}
  \Delta C = 45 \times 10^9 \times \$1.609 = \$72.4 \text{ billion/year (direct cost reduction)}
\end{equation}

The \$135 billion headline figure reconciles with this calculation by including the revenue recovery on previously empty miles: trucks that were running empty and are now carrying loads generate revenue. If the loaded freight generates \$1.40 per loaded mile (the approximate industry-average revenue per loaded truck-mile in 2023 \citep{DAT_spotrates2023}), then 45 billion newly loaded miles generate:

\begin{equation}
  \Delta R = 45 \times 10^9 \times \$1.40 = \$63.0 \text{ billion/year (revenue recovery)}
\end{equation}

Combined: \$72.4 billion in cost avoidance plus \$63.0 billion in revenue recovery equals approximately \$135 billion in aggregate freight efficiency gain---the headline figure \citep{ROUTE_C1}.

\subsection{Corridor-Level Analysis}

The efficiency gain is not uniformly distributed across corridors. Corridors with the highest structural flow imbalance and the highest truck volumes generate the largest absolute savings. Table~\ref{tab:corridors} ranks the top 10 T1/T2 corridors by estimated annual efficiency gain from relay hub pre-matching.

\begin{table}[h]
\centering
\caption{Estimated Annual Efficiency Gain by Corridor from Relay Hub Pre-Matching}
\label{tab:corridors}
\begin{tabular}{lccccc}
\toprule
Corridor & Daily Trucks & Empty Rate & Target & Gain (pp) & Annual Gain (\$B) \\
\midrule
I-80 (Chicago--SF) & 28,000 & 34\% & 19\% & 15 & \$11.2 \\
I-95 (Boston--Miami) & 31,000 & 31\% & 18\% & 13 & \$10.4 \\
I-10 (LA--Houston) & 22,000 & 38\% & 21\% & 17 & \$8.7 \\
I-40 (LA--Charlotte) & 18,000 & 35\% & 20\% & 15 & \$6.6 \\
I-29 (Fargo--KC) & 8,500 & 47\% & 28\% & 19 & \$3.9 \\
I-35 (Laredo--KC) & 16,000 & 33\% & 19\% & 14 & \$5.5 \\
I-75 (Atlanta--Detroit) & 19,000 & 29\% & 17\% & 12 & \$5.6 \\
I-55 (Memphis--Chicago) & 12,000 & 32\% & 19\% & 13 & \$3.8 \\
I-90 (Seattle--Chicago) & 10,000 & 36\% & 20\% & 16 & \$3.9 \\
I-70 (Baltimore--KC) & 11,000 & 31\% & 19\% & 12 & \$3.2 \\
\midrule
Top 10 Total & & & & & \$62.8 \\
\bottomrule
\end{tabular}
\end{table}

\noindent Note: Annual gains computed as daily-trucks $\times$ (gain in pp / 100) $\times$ 365 $\times$ average one-way haul distance (800 miles) $\times$ (\$1.609/mile + \$1.40/mile). Figures are estimates; see Section~\ref{sec:methods-note} for data sources and uncertainty discussion.

The top 10 corridors alone generate approximately \$63 billion of the \$135 billion aggregate gain. The remaining \$72 billion accrues on the long tail of regional and secondary corridors that collectively account for half of all truck-miles.

\subsection{Who Captures the Surplus}

In competitive freight markets, efficiency gains from load matching are competed away by shippers (through rate pressure) and drivers (through higher acceptance rates on previously unacceptable loads). The distributional analysis is:

\textbf{Carriers}: Receive direct cost reduction from fewer empty miles (\$72.4 billion in avoided operating cost). Competitive pressure from other carriers adopting the same platform will push approximately 60--70\% of this benefit to shippers through rate competition, leaving carriers with approximately \$22--29 billion in permanent margin improvement.

\textbf{Shippers}: Benefit from lower rates as carriers share efficiency gains. Estimated benefit: \$43--50 billion per year in lower freight costs. This benefit flows through to consumer prices in proportion to freight cost share by sector.

\textbf{Drivers}: Relay hub pre-matching reduces driver search time and improves home-time predictability. Drivers currently spending 30--60 minutes per load searching for return freight recover approximately 20 minutes of productive time per load event. At \$89/hr driver wage \citep{ATRI_costs2024} and approximately 200 load events per driver per year, the per-driver time recovery is approximately \$593/year. Multiplied by 3.5 million long-haul drivers, the aggregate driver time recovery is approximately \$2 billion/year---a small share of the total but meaningful at the individual level.

\subsection{Comparison to Managed-Lane Capital Cost}

The Interstate 2.0 managed-lane network is estimated to require approximately \$121 billion in capital investment across the top T1/T2 corridors \citep{ROUTE_C1}. The relay hub load-matching function generates \$113--\$135 billion per year in efficiency gains---exceeding the one-time capital cost in the first year of full deployment. The managed-lane investment is justified on reliability and throughput grounds; the load-matching function provides a recurring economic return that makes the infrastructure investment double-purpose.

This comparison understates the relay hub's load-matching value relative to its cost because the matching platform requires only software and data infrastructure at the hub---not additional construction. The relay terminal itself is built for driver-relay and intermodal purposes; the load-matching platform is an incremental cost of approximately \$2--5 million per hub in software and communications infrastructure.
